% Unit 065 terjemahan Bahasa Indonesia dari chapter5.tex baris 141--349.
% Sumber: Methods of Algebra, Volume 2, commit
% 9a5803ff77dd3257484cb177f851a73770a59dd3 (CC BY 4.0).
% stable-unit-id: o014.aljabr2.chapter5.general-spectral-sequences
% Cakupan: Bagian 5.2 lengkap.

% segment-id: o014.li2.chapter5.general-spectral-sequences.h001
\section{Definisi umum barisan spektral}\label{sec:spectral-sequence}
\hypertarget{o014.aljabr2.chapter5.general-spectral-sequences}{ }

% segment-id: o014.li2.chapter5.general-spectral-sequences.p001
Kita mulai dengan definisi umum objek diferensial. Definisi ini melibatkan
kategori aditif dengan translasi $(\mathcal{A},T)$ sebagaimana dalam
Definisi~\ref{def:category-with-translation}, serta morfisme berderajat
sebagaimana dalam Definisi~\ref{def:morphism-with-degree}.

% segment-id: o014.li2.chapter5.general-spectral-sequences.d001
\begin{definition}\label{def:differential-object}
	\index{objek diferensial (differential object)}
	\index[sym1]{ATd@$\mathcal{A}_d$, $(\mathcal{A}, T)_d$}
	Sebuah \emph{objek diferensial} pada kategori aditif dengan translasi
	$(\mathcal{A},T)$ adalah data $(X,d)$, dengan $X\in\Obj(\mathcal{A})$ dan
	$d:X\xrightarrow{+1}X$ yang memenuhi $d^2=0$.
	\begin{itemize}
		\item Morfisme dari $(X,d)$ ke $(X',d')$ adalah morfisme
		$\alpha:X\to X'$ yang memenuhi $d'\alpha=\alpha d$.
		\item Semua objek diferensial membentuk kategori $(\mathcal{A},T)_d$;
		kita juga menulis
		$\mathcal{A}_d:=(\mathcal{A},\identity_{\mathcal{A}})_d$.
	\end{itemize}
\end{definition}

% segment-id: o014.li2.chapter5.general-spectral-sequences.p002
Jika $\mathcal{A}$ adalah kategori-$\cate{Ab}$ (atau
kategori-$\Bbbk\dcate{Mod}$, dengan $\Bbbk$ gelanggang komutatif), maka
$(\mathcal{A},T)_d$ juga demikian. Funktor pelupa
$(\mathcal{A},T)_d\to\mathcal{A}$ memetakan $(X,d)$ ke $X$.

% segment-id: o014.li2.chapter5.general-spectral-sequences.t001
\begin{proposition}
	\index[sym1]{HXd@$\Hm(X,d)$}
	Funktor pelupa $(\mathcal{A},T)_d\to\mathcal{A}$ menciptakan semua
	$\varinjlim$ dan $\varprojlim$ (Definisi~\ref{def:create-limit}).
\end{proposition}
\begin{proof}
	Tinjau kembali bukti Lema~\ref{prop:limits-cplx}.
\end{proof}

% segment-id: o014.li2.chapter5.general-spectral-sequences.p003
Karena itu, jika $\mathcal{A}$ adalah kategori abelian, maka
$(\mathcal{A},T)_d$ juga kategori abelian, dan
$(\mathcal{A},T)_d\to\mathcal{A}$ eksak.

% segment-id: o014.li2.chapter5.general-spectral-sequences.p004
Dalam kasus kategori abelian, salah satu ciri objek diferensial adalah bahwa
objek tersebut memiliki homologi atau kohomologi.

% segment-id: o014.li2.chapter5.general-spectral-sequences.d002
\begin{definition}
	Misalkan $\mathcal{A}$ kategori abelian. Kita mendefinisikan funktor aditif
	\[ \Hm:(\mathcal{A},T)_d\to\mathcal{A},\qquad
	   \Hm(X,d):=\Ker(d)/\Image(T^{-1}d). \]
\end{definition}

% segment-id: o014.li2.chapter5.general-spectral-sequences.p005
Kita menyebut diagram segitiga dalam kategori abelian $\mathcal{A}$
$\begin{tikzcd}[column sep=small, row sep=small]
	A \arrow[rr] & & B \arrow[ld] \\
	& C \arrow[lu, "+m"] &
\end{tikzcd}$
\emph{eksak} apabila
\[ \cdots T^{k-m}C\to T^kA\to T^kB\to T^kC\to T^{k+m}A\cdots \]
merupakan barisan eksak yang memanjang tak hingga ke kedua arah. Diagram
segitiga berderajat yang lebih umum diperlakukan dengan cara yang sama.
\index{barisan eksak}

% segment-id: o014.li2.chapter5.general-spectral-sequences.l001
\begin{lemma}\label{prop:diff-obj-triangle}
	Misalkan $(\mathcal{A},T)$ suatu kategori abelian dengan translasi. Jika
	$0\to(X',d')\to(X,d)\to(X'',d'')\to0$ adalah barisan eksak pendek dalam
	$(\mathcal{A},T)_d$, maka diagram berikut eksak pada setiap titik sudut:
	\[\begin{tikzcd}[column sep=tiny]
		\Hm(X',d') \arrow[rr] & & \Hm(X,d) \arrow[ld] . \\
		& \Hm(X'',d'') \arrow[lu, "{+1}"] &
	\end{tikzcd}\]
\end{lemma}
\begin{proof}
	Dalam $\mathcal{A}$, perhatikan barisan kompleks
	\[ 0\to(T^nX',T^nd')_n\to(T^nX,T^nd)_n
	   \to(T^nX'',T^nd'')_n\to0. \]
	Barisan ini eksak derajat demi derajat, sehingga merupakan barisan eksak
	pendek. Terapkan Proposisi~\ref{prop:long-exact-sequence-ses} kepadanya.
\end{proof}

% segment-id: o014.li2.chapter5.general-spectral-sequences.e001
\begin{example}[Kompleks sebagai objek diferensial]\label{eg:cplx-as-diff}
	\index{objek bergradasi diferensial}
	Untuk kategori aditif $\mathcal{A}$, perhatikan $\mathcal{A}^{\Z}$ bersama
	funktor translasi $T:(X^n)_n\mapsto(X^{n+1})_n$. Berdasarkan uraian dalam
	\S\ref{sec:additive-cplx}, objek dalam $(\mathcal{A}^{\Z},T)_d$ juga
	disebut objek bergradasi diferensial, dan terdapat isomorfisme kategori
	$\cate{C}(\mathcal{A})\simeq(\mathcal{A}^{\Z},T)_d$. Jika $\mathcal{A}$
	adalah kategori abelian, jelas bahwa
	\[ \Hm(X,d)=\left(\Hm^p(X)\right)_{p\in\Z}. \]
\end{example}

% segment-id: o014.li2.chapter5.general-spectral-sequences.p006
Sekarang kita dapat memberikan definisi barisan spektral. Untuk menyederhanakan
pembahasan, mula-mula kita tinjau kasus tanpa translasi, atau setara dengan itu,
kasus morfisme tanpa derajat; dengan kata lain, kita ambil
$T=\identity$.

% segment-id: o014.li2.chapter5.general-spectral-sequences.d003
\begin{definition}[J.\ Leray, J.-L.\ Koszul]\label{def:SS-abstract}
	\index{barisan spektral (spectral sequence)}
	Misalkan $\mathcal{A}$ kategori abelian dan $a\in\Z$. Sebuah
	\emph{barisan spektral} yang dimulai dari $a$ adalah data
	$\mathscr{E}=(E_r,d_r)_{r\in\Z_{\geq a}}$ bersama $(t_r)_{r\geq a+1}$,
	dengan
	\begin{compactitem}
		\item $(E_r,d_r)\in\Obj(\mathcal{A}_d)$;
		\item $t_r:\Hm(E_{r-1},d_{r-1})\rightiso E_r$ untuk $r\geq a+1$.
	\end{compactitem}
	Data $a$ dan $(t_r)_r$ sering dihilangkan dari notasi.

	Morfisme barisan spektral $\mathscr{E}\to\mathscr{E}'$ adalah keluarga
	morfisme dalam $\mathcal{A}_d$
	$\varphi_r:(E_r,d_r)\to(E'_r,d'_r)$ yang memenuhi
	$t_r\Hm(\varphi_{r-1})=\varphi_rt_r$ untuk $r\geq a+1$.
\end{definition}

% segment-id: o014.li2.chapter5.general-spectral-sequences.p007
Pasangan $(E_r,d_r)$ lazim disebut \emph{halaman ke-$r$} dari barisan spektral
$\mathscr{E}$; beralih dari $E_r$ ke $E_{r+1}$ berarti membalik halaman.
Secara abstrak, halaman tempat barisan spektral mulai tidaklah penting, tetapi
dalam penerapan pilihan tersebut bergantung pada konteksnya.

% segment-id: o014.li2.chapter5.general-spectral-sequences.p008
Jika $d_r$ dalam barisan spektral dipandang sebagai morfisme
$\Hm(E_{r-1},d_{r-1})\to\Hm(E_{r-1},d_{r-1})$, maka
$\Ker(d_r)\supset\Image(d_r)$ masing-masing bersesuaian dengan subobjek dari
$\Ker(d_{r-1})$ yang memuat $\Image(d_{r-1})$. Demi kemudahan,
anggap barisan spektral dimulai dari $a=0$. Definisikan $Z_1$ (berturut-turut
$B_1$) sebagai $\Ker(d_0)$ (berturut-turut $\Image(d_0)$), lalu definisikan
$Z_2$ (berturut-turut $B_2$) sebagai prapeta $\Ker(d_1)$ (berturut-turut
$\Image(d_1)$) terhadap $E_0\supset Z_1\twoheadrightarrow E_1$, dan seterusnya.
Iterasi ini memberikan
\begin{equation}\label{eqn:spectral-sequence-Z-B}\begin{gathered}
	0=:B_0\subset B_1\subset B_2\subset\cdots\subset Z_2\subset Z_1
	\subset Z_0=:E_0, \\
	Z_r\leftrightarrow\Ker(d_{r-1}),\quad
	B_r\leftrightarrow\Image(d_{r-1}),\quad Z_r/B_r\simeq E_r.
\end{gathered}\end{equation}
\index[sym1]{ZBE@$Z_r$, $B_r$, $E_r$}

% segment-id: o014.li2.chapter5.general-spectral-sequences.p009
\begin{itemize}
	\item Jika $Z_\infty:=\bigcap_rZ_r$ dan $B_\infty:=\bigcup_rB_r$ ada,
	\emph{limit} barisan spektral $\mathscr{E}$ didefinisikan sebagai
	$E_\infty:=Z_\infty/B_\infty$.
	\item Jika ada $r$ sedemikian sehingga $r'\geq r\implies d_{r'}=0$, maka
	$\mathscr{E}$ dikatakan \emph{mengalami degenerasi pada halaman $E_r$}.
	Dalam keadaan ini jelas bahwa $Z_r=Z_{r+1}=\cdots$ dan
	$B_r=B_{r+1}=\cdots$, sehingga $Z_r=Z_\infty$ dan $B_r=B_\infty$;
	sebagai akibatnya, $E_r=E_\infty$.
	\index{barisan spektral!degenerasi (degenerate)}
\end{itemize}

% segment-id: o014.li2.chapter5.general-spectral-sequences.p010
Ketika membahas barisan spektral bergradasi homologis, kelak kita akan memakai
notasi $(E^r,d^r)_r$ serta $B^r,Z^r$.

% segment-id: o014.li2.chapter5.general-spectral-sequences.t002
\begin{proposition}\label{prop:ss-isom-infty}
	Misalkan $\varphi_r:\mathscr{E}\to\mathscr{E}'$ adalah morfisme barisan
	spektral. Jika $\varphi_r$ merupakan isomorfisme, maka $\varphi_s$ juga
	isomorfisme untuk $s\geq r$; jika limitnya ada, maka
	$\varphi_\infty:E_\infty\to E'_\infty$ juga isomorfisme.
\end{proposition}
\begin{proof}
	Isomorfisme $\varphi_r$ menginduksi isomorfisme
	$\Hm(E_r,d_r)\to\Hm(E'_r,d'_r)$, yang tepat merupakan
	$\varphi_{r+1}$. Tanpa mengurangi keumuman, anggap barisan spektral dimulai
	dari $r$. Maka $(\varphi_s)_{s\geq r}$ adalah isomorfisme barisan spektral,
	dan pernyataan mengenai $\varphi_\infty$ langsung mengikuti.
\end{proof}

% segment-id: o014.li2.chapter5.general-spectral-sequences.p011
Lebih umum, untuk kategori abelian dengan translasi $(\mathcal{A},T)$,
definisi barisan spektral dapat diperluas ke keadaan
$(E_r,d_r)\in\Obj((\mathcal{A},T^{a_r})_d)$, dengan
$a_1,a_2,\ldots$ suatu barisan bilangan bulat. Definisi $\Hm(E_r,d_r)$ serta
$B_r,Z_r$ di atas tidak berubah.

% segment-id: o014.li2.chapter5.general-spectral-sequences.p012
Dalam perluasan yang akan kita perlukan, kita bahkan dapat membiarkan
$\mathcal{A}$ dilengkapi keluarga autoisomorfisme saling komutatif
$T_1,\ldots,T_n$, dan
\[ d_r:E_r\xrightarrow{+\vec{a}_r}E_r,\qquad
	\vec{a}_r=(a_{r,1},\ldots,a_{r,n})\in\Z^n. \]
Mengikuti gagasan ini, kita memperkenalkan dua versi yang paling umum dipakai;
keduanya melibatkan struktur bergradasi.

% segment-id: o014.li2.chapter5.general-spectral-sequences.d004
\begin{definition}[Barisan spektral: versi bergradasi dan bergradasi ganda]
\label{def:graded-spectral-sequence}
	\index{barisan spektral!bergradasi, bergradasi ganda (graded, bigraded)}
	\index[sym1]{drpq@$d_r^{p,q}$, $d^r_{p,q}$}
	Misalkan $\mathcal{A}$ kategori abelian. Barisan spektral bergradasi
	kohomologis dan barisan spektral bergradasi ganda kohomologis, keduanya
	ditulis sebagai $\mathscr{E}=(E_r,d_r)_r$, didefinisikan sebagai berikut.
	\[\begin{array}{|c|c|c|c|} \hline
		\text{Versi} & E_r & d_r & \text{Derajat morfisme} \\ \hline
		\text{Bergradasi} & (E_r^p)_{p\in\Z}\in\Obj(\mathcal{A}^{\Z})
		& (d_r^p)_p & d_r^p:E_r^p\xrightarrow{+r}E_r^{p+r} \\
		\text{Bergradasi ganda} &
		(E_r^{p,q})_{(p,q)\in\Z^2}\in\Obj(\mathcal{A}^{\Z\times\Z})
		& (d_r^{p,q})_{(p,q)} &
		d_r^{p,q}:E_r^{p,q}\xrightarrow{+(r,-r+1)}E_r^{p+r,q-r+1}
		\\ \hline
	\end{array}\]
	Dalam kedua kasus disyaratkan $(d_r)^2=0$, dengan komposisi dipahami
	sebagai komposisi morfisme berderajat, dan datanya memuat isomorfisme
	tertentu $\Hm(E_{r-1},d_{r-1})\rightiso E_r$.

	Secara dual, barisan spektral bergradasi (atau bergradasi ganda) homologis
	didefinisikan sebagai data $(E^r,d^r)_{r\geq1}$ berikut, dengan syarat yang
	sama bahwa $(d^r)^2=0$ dan dengan isomorfisme tertentu
	$\Hm(E^r,d^r)\simeq E^{r+1}$:
	\[\begin{array}{|c|c|c|c|} \hline
		\text{Versi} & E^r & d^r & \text{Derajat morfisme} \\ \hline
		\text{Bergradasi} & (E^r_p)_{p\in\Z}\in\Obj(\mathcal{A}^{\Z})
		& (d^r_p)_p & d^r_p:E^r_p\xrightarrow{-r}E^r_{p-r} \\
		\text{Bergradasi ganda} &
		(E^r_{p,q})_{(p,q)\in\Z^2}\in\Obj(\mathcal{A}^{\Z\times\Z})
		& (d^r_{p,q})_{(p,q)} &
		d^r_{p,q}:E^r_{p,q}\xrightarrow{+(-r,r-1)}E^r_{p-r,q+r-1}
		\\ \hline
	\end{array}\]
	Morfisme antarbarisan spektral ini didefinisikan dengan cara
	biasa dan harus kompatibel dengan isomorfisme yang ditentukan.
\end{definition}

% segment-id: o014.li2.chapter5.general-spectral-sequences.p013
Jika digambar dalam koordinat $(p,q)$, arah $d_r$ dan $d^r$ pada barisan
spektral bergradasi ganda adalah sebagai berikut.
\begin{center}\begin{tikzpicture}[>=Latex, scale=0.7, baseline=(O)]
	\draw[very thin, gray!50] (-1,1) grid (4,-2);
	\draw[thick, ->] (0,0) -- (0,1) node[left] {$d_0$};
	\draw[thick, ->] (0,0) -- (1,0) node[right] {$d_1$};
	\draw[thick, ->] (0,0) -- (2,-1) node[right] {$d_2$};
	\draw[thick, ->] (0,0) -- (3,-2) node[right] {$d_3$};
	\node at (1.5,-2.5) {$E_r^{p,q}$};
	\coordinate (O) at (1.5,-0.5);
\end{tikzpicture} \qquad \begin{tikzpicture}[>=Latex, scale=0.7, baseline=(O)]
	\draw[very thin, gray!50] (1,-1) grid (-4,2);
	\draw[thick, ->] (0,0) -- (0,-1) node[right] {$d^0$};
	\draw[thick, ->] (0,0) -- (-1,0) node[left] {$d^1$};
	\draw[thick, ->] (0,0) -- (-2,1) node[left] {$d^2$};
	\draw[thick, ->] (0,0) -- (-3,2) node[left] {$d^3$};
	\node at (-1.5,-1.5) {$E^r_{p,q}$};
	\coordinate (O) at (-1.5,0.5);
\end{tikzpicture}\end{center}

% segment-id: o014.li2.chapter5.general-spectral-sequences.p014
Dalam versi mana pun, mengikuti pola sebelumnya kita dapat mendefinisikan
objek-objek $B_r^p\subset Z_r^p$, $B_r^{p,q}\subset Z_r^{p,q}$,
$E_\infty^p$, $E_\infty^{p,q}$, serta
$B^r_p\subset Z^r_p$, $B^r_{p,q}\subset Z^r_{p,q}$,
$E^\infty_p$, $E^\infty_{p,q}$, dan seterusnya.

% segment-id: o014.li2.chapter5.general-spectral-sequences.e002
\begin{example}[Kasus dengan lebar hingga]\label{eg:SS-finite-width}
	Dalam banyak situasi umum, suku-suku taknol suatu barisan spektral
	bergradasi ganda terkonsentrasi pada pita horizontal (atau vertikal)
	selebar $s$, dengan $s\in\Z_{\geq0}$, seperti pada gambar berikut:
	\[\begin{tikzpicture}[baseline]
		\matrix (M) [matrix of math nodes, right delimiter=\}] {
			\cdots & \bullet & \bullet & \bullet & \cdots \\
			\vdots & \vdots & \vdots & \vdots & \vdots \\
			\cdots & \bullet & \bullet & \bullet & \cdots \\
		};
		\node[right=2em] at (M-2-5) {\scriptsize \text{$s$ baris}};
	\end{tikzpicture} \quad \text{atau} \quad \begin{tikzpicture}[baseline]
		\matrix (M) [matrix of math nodes, below delimiter=\}] {
			\vdots & \vdots & \vdots \\
			\bullet & \cdots & \bullet \\
			\bullet & \cdots & \bullet \\
			\vdots & \vdots & \vdots \\
		};
		\node[below=2em] at (M-4-2) {\scriptsize \text{$s$ kolom}};
	\end{tikzpicture}\]
	Dengan memperhatikan arah panah, tampak bahwa syarat $r>s$ (atau $r\geq s$)
	mengakibatkan $d_r^{p,q}=0$ untuk semua $p,q$; dalam keadaan ini barisan
	spektral mengalami degenerasi pada halaman $E_r$. Kasus homologis serupa.
\end{example}

% segment-id: o014.li2.chapter5.general-spectral-sequences.d005
\begin{definition}\label{def:bdd-ss}
	\index{barisan spektral!terbatas (bounded)}
	Misalkan $\mathscr{E}$ adalah barisan spektral bergradasi ganda kohomologis
	(atau homologis), dan $r\in\Z$. Jika untuk setiap $n\in\Z$ hanya terdapat
	berhingga banyak pasangan $(p,q)\in\Z^2$ yang memenuhi $p+q=n$ dan
	$E_r^{p,q}\neq0$ (atau $E^r_{p,q}\neq0$), maka $E_r$ (atau $E^r$)
	disebut \emph{terbatas}.
\end{definition}

% segment-id: o014.li2.chapter5.general-spectral-sequences.p015
Dari $E_{r+1}\simeq\Hm(E_r,d_r)$ terlihat bahwa
$E_r^{p,q}=0\implies E_{r+1}^{p,q}=0$. Khususnya, keterbatasan $E_r$
mengakibatkan keterbatasan $E_{r+1}$. Kasus homologis serupa.

% segment-id: o014.li2.chapter5.general-spectral-sequences.t003
\begin{proposition}
	Misalkan barisan spektral bergradasi ganda kohomologis (atau homologis)
	$\mathscr{E}$ mempunyai $E_r$ (atau $E^r$) yang terbatas. Untuk setiap
	$(p,q)\in\Z^2$ terdapat $r(p,q)$ sedemikian sehingga, jika
	$r\geq r(p,q)$, maka $E_r^{p,q}=E_{r+1}^{p,q}$ (atau
	$E^r_{p,q}=E^{r+1}_{p,q}$). Akibatnya, limit $\mathscr{E}$ ada dan
	memenuhi $E_r^{p,q}=E_\infty^{p,q}$ (atau
	$E^r_{p,q}=E^\infty_{p,q}$).
\end{proposition}
\begin{proof}
	Cukup bahas kasus $E_r$. Misalkan $n$ tetap. Karena sifat keterbatasan, untuk
	$r$ yang cukup besar dan semua $(p,q)$ dengan $p+q=n$, berlaku
	$E_r^{p-r,q+r-1}=0$, sehingga $B_r^{p,q}=0$. Dengan cara yang sama, untuk
	$r$ yang cukup besar berlaku $E_r^{p+r,q-r+1}=0$, sehingga
	$Z_r^{p,q}=E_r^{p,q}$. Ini membuktikan pernyataan tersebut.
\end{proof}

% segment-id: o014.li2.chapter5.general-spectral-sequences.d006
\begin{definition}\label{def:quadrant-ss}
	Misalkan $\mathscr{E}$ adalah barisan spektral bergradasi ganda kohomologis
	(atau homologis), dan $r\in\Z$. Jika $E_r^{p,q}\neq0$ (atau
	$E^r_{p,q}\neq0$) hanya mungkin untuk $p,q\geq0$, maka $E_r$ (atau $E^r$)
	dikatakan terletak di kuadran pertama. Kasus untuk kuadran-kuadran lain
	didefinisikan dengan cara serupa.
\end{definition}

% segment-id: o014.li2.chapter5.general-spectral-sequences.p016
$E_r$ yang terletak di kuadran pertama atau ketiga jelas terbatas. Dari
$E_{r+1}\simeq\Hm(E_r,d_r)$ tampak bahwa jika $E_r$ terletak di kuadran
pertama, dan seterusnya, maka $E_{r+1}$ juga demikian. Kasus homologis serupa.

% segment-id: o014.li2.chapter5.general-spectral-sequences.r001
\begin{remark}[Perhitungan tepi]\label{rem:edge-computing}
	Barisan spektral yang terletak di kuadran pertama sangat umum dijumpai.
	Untuk versi kohomologis, suku tepinya $E_r^{\bullet,0}$ dan
	$E_r^{0,\bullet}$ mempunyai sifat khusus. Misalkan $p,q\geq1$. Dengan
	menelaah arah $d_r$ dan mengingat bahwa
	$E_{r+1}\simeq\Hm(E_r,d_r)$, kita memperoleh
	\begin{equation}\label{eqn:edge-ss-coh}\begin{array}{cc}
		E_2^{p,0}\twoheadrightarrow E_3^{p,0}\twoheadrightarrow\cdots
		\twoheadrightarrow E_{p+1}^{p,0}=E_\infty^{p,0}
		& \text{($\because$ semua $d$ yang keluar bernilai $0$)}, \\
		E_\infty^{0,q}=E_{q+2}^{0,q}\hookrightarrow E_{q+1}^{0,q}
		\hookrightarrow\cdots\hookrightarrow E_2^{0,q}\hookrightarrow E_1^{0,q}
		& \text{($\because$ semua $d$ yang masuk bernilai $0$)}.
	\end{array}\end{equation}
	Barisan spektral bergradasi ganda homologis di kuadran pertama memiliki
	sifat yang bersesuaian:
	\begin{equation}\label{eqn:edge-ss-ho}\begin{array}{cc}
		E^\infty_{p,0}=E^{p+1}_{p,0}\hookrightarrow E^p_{p,0}
		\hookrightarrow\cdots\hookrightarrow E^3_{p,0}\hookrightarrow E^2_{p,0}
		& \text{($\because$ semua $d$ yang masuk bernilai $0$)}, \\
		E^1_{0,q}\twoheadrightarrow E^2_{0,q}\twoheadrightarrow\cdots
		\twoheadrightarrow E^{q+2}_{0,q}=E^\infty_{0,q}
		& \text{($\because$ semua $d$ yang keluar bernilai $0$)}.
	\end{array}\end{equation}
	\index{morfisme tepi (edge morphism)}

	Morfisme yang berasal dari \eqref{eqn:edge-ss-coh} atau
	\eqref{eqn:edge-ss-ho} secara bersama-sama disebut \emph{morfisme tepi}.
	Dari uraian di atas diperoleh barisan eksak berikut ($p\geq2$):
	\begin{equation}\label{eqn:low-exact-coh}
		0\to E^{0,p-1}_\infty\to E^{0,p-1}_p\xrightarrow{d}E^{p,0}_p
		\to E^{p,0}_\infty\to0,
	\end{equation}
	dengan semua morfisme selain $d$ merupakan morfisme tepi. Terdapat pula
	versi homologis dari barisan eksak tersebut ($p\geq2$):
	\begin{equation}\label{eqn:low-exact-ho}
		0\to E^\infty_{p,0}\to E^p_{p,0}\xrightarrow{d}E^p_{0,p-1}
		\to E^\infty_{0,p-1}\to0.
	\end{equation}
	Pembaca pemula sangat dianjurkan untuk menggambar sendiri arah morfisme
	$d$ yang melibatkan suku-suku ini, menentukan kapan morfisme tersebut
	bernilai $0$, lalu memverifikasi semua pernyataan di atas.
\end{remark}
